\documentclass[11pt]{article}

\usepackage[review]{acl}
\usepackage{times}
\usepackage{latexsym}
\usepackage[T1]{fontenc}
\usepackage[utf8]{inputenc}
\usepackage{microtype}
\usepackage{inconsolata}
\usepackage{graphicx}
\usepackage{booktabs}
\usepackage{array}
\usepackage{amsmath}
\usepackage{amssymb}

\title{Where Investment Decision Steering Is Accessible: Layer-Localized Stance Axes and Multi-dimensional Operators in an LLM}

\author{First Author \\
  Affiliation / Address line 1 \\
  Affiliation / Address line 2 \\
  \texttt{email@domain} \\
  \And
  Second Author \\
  Affiliation / Address line 1 \\
  Affiliation / Address line 2 \\
  \texttt{email@domain} \\
}

\begin{document}
\maketitle

\begin{abstract}
Changing only the named company can alter an LLM's investment decision even when the explicit
financial evidence is held fixed. We study how this decision stance can be accessed and controlled
at inference time. Using a fixed JSON decision protocol, we first localize the layer and prompt
span where residual interventions have leverage over a fixed buy/sell readout. We then compare a
single-neuron intervention, a one-dimensional difference-in-means direction, and a four-dimensional
contrastive concept cone. Throughout, we distinguish movement in the fixed-token buy/sell margin
from changes in the complete greedy decision. The study is designed to characterize the accessible
control surface of investment decisions, rather than to identify the origin of entity bias or to
claim entity-selective debiasing.
\end{abstract}

\section{Introduction}

Language models are increasingly used to summarize evidence, compare alternatives, and recommend
actions. In an investment prompt, changing only the named company can change the model's buy/sell
output even when the explicit evidence is held fixed. This entity-dependent behavior motivates a
more focused question: how can an LLM's investment stance be controlled at inference time, and
where in its residual representation does such control have leverage?

Prior work showed that a scalar intervention on one neuron can increase or decrease its activation
and shift a model's aggregate buy/sell stance; Park et al. (2026) refer to this control as an
investment dial. This result demonstrates that investment stance is controllable without changing
the prompt or model weights. It does not establish which layer and prompt span provide a useful
intervention site, whether a multi-dimensional operator offers a broader control surface than a
single neuron, or whether a change in a fixed-token readout reliably changes the model's generated
decision.

We are motivated by the gap between observing that an LLM's investment stance is controllable in
its internal representation and knowing whether that control can be localized, extended beyond a
single neuron, and translated into generated decisions. The financial setting adds a further
requirement: a changed recommendation should be considered alongside its sensitivity to the evidence
provided in the prompt. A neuron-level intervention can reveal a control channel without showing
which evidence supports the resulting output. We therefore evaluate generated decisions and evidence
sensitivity alongside internal margin changes.

Recent refusal-steering work provides a methodological precedent for this question. Arditi et al.
(2024) showed that adding or removing a residual-stream refusal direction can bidirectionally
control generated behavior, while Wollschl\"ager et al. (2025) argued that refusal may occupy a
multi-dimensional concept cone rather than a single functional direction. We transfer this
intervention perspective from safety refusal to investment stance. The transfer is non-trivial:
refusal is a safety behavior elicited by harmful or benign requests, whereas our target is a
template-conditioned buy/sell generation. Our operators are constructed from company-level
decision-margin contrasts, and we evaluate both fixed-token readouts and complete greedy decisions.

We study inference-time representation steering under a fixed JSON decision template. We first use
layer-wise residual patching to identify where differences associated with the entity, evidence,
instruction, and final spans transfer a fixed buy/sell readout. We then compare three steering
mechanisms at the selected layer: the single-neuron scalar intervention described above; a token-wise
difference-in-means (DIM) direction formed from low- and high-margin companies; and a four-dimensional
cone formed from sector-demeaned contrastive-SVD directions.

This paper makes three contributions:

\begin{enumerate}
    \item We provide a layer- and span-localized analysis of where residual interventions have
    leverage over a fixed-token investment readout.
    \item We compare a single-neuron scalar intervention, a one-dimensional DIM direction, and a
    multi-dimensional concept cone under a common steering protocol and matched-norm controls.
    \item We separate fixed-token margin movement from generated decision changes, documenting the
    boundary at which internal readout changes do or do not produce greedy JSON decision flips.
\end{enumerate}

Taken together, we frame investment stance as an inference-time steering problem. Entity-dependent
behavior motivates the setting, while the main focus is where decision steering is accessible, how
different operators control it, and how reliably internal readout changes translate into generated
decisions.

\section{Related Work}

\subsection{Entity-Conditioned Investment Decisions}

A growing body of work documents that language-model outputs can depend on entity identity even when
the surrounding task is held fixed. In financial settings, this motivates careful separation between
entity-conditioned behavior and the internal mechanisms used to control a decision. We use this
phenomenon as the motivation for the steering problem; we do not treat the present study as a
localization or debiasing account of entity-specific factual memory.

\subsection{Inference-Time Representation Steering}

Our intervention design is related to methods that modify hidden representations during inference
without updating model weights. Refusal removal and refusal steering provide a methodological
precedent for testing whether a behavioral tendency can be controlled through a residual-space
direction. Other work studies whether a concept is adequately represented by a single direction or
requires a multi-dimensional subspace. We transfer these ideas to investment stance and evaluate
both internal readout movement and complete generated decisions.

\subsection{Layer-Wise Activation Patching}

Layer-wise activation patching and causal tracing provide tools for identifying where information
that affects an output is transferred through a transformer. We use this family of methods to select
an intervention site before comparing steering operators. The selected site is interpreted as a
location with measured steering leverage, not as proof that the site is the origin of
entity-dependent behavior.

\section{Study Design}

We formulate investment stance steering as an intervention on the residual representation during
prompt prefill. The current instantiation uses Qwen3.5-4B, a 32-layer decoder model with hidden size
2560. The protocol selects the intervention layer separately for each model; the layer index reported
for Qwen3.5-4B should therefore be read as a model-specific result.

\subsection{Task and Measurements}

Each prompt contains a company identity, a financial evidence block, and an instruction requiring a
single JSON object with a \texttt{buy} or \texttt{sell} decision and a short reason. We measure a
continuation margin at the answer prefix,

\begin{equation}
M = \log p(\text{buy}) - \log p(\text{sell}),
\end{equation}

using the same fixed answer-token scorer for every condition. This margin measures movement in the
fixed-token readout. We separately generate the complete response with greedy decoding ($T=0$, at
most 48 new tokens), parse the JSON decision, and record whether the output is \texttt{buy},
\texttt{sell}, or unparsed. A positive margin and a generated \texttt{buy} are therefore separate
outcomes.

The canonical balanced prompt uses the frozen decision template. Evidence-condition probes vary the
evidence block while keeping the company and decision format fixed. All pooled scenario results must
use a single frozen prompt renderer and a recorded prompt hash.

\subsection{Cohort and Operator Construction}

Operator construction starts from the 200-company cohort in
\texttt{entity-to-dial-heldout-transfer-v1-01}. Companies are ranked by their clean continuation
margin. The token-wise DIM direction uses the ten highest- and ten lowest-margin companies. At each
of the 100 instruction-span token positions, we subtract the mean residual state of the low-margin
group from the mean residual state of the high-margin group:

\begin{equation}
v_{\mathrm{DIM}}[p] = \mu_{\mathrm{high}}[p] - \mu_{\mathrm{low}}[p].
\end{equation}

The multi-dimensional operator uses the twenty highest- and twenty lowest-margin companies. We
subtract the mean state of each sector before forming paired high-minus-low residual differences,
perform SVD independently at each instruction token, retain the leading four directions, and align
their signs with the positive contrast. The cone centroid is the normalized sum of these four
directions.

The current construction and pilot evaluation draw from the same 200-company cohort. They therefore
measure steering feasibility on the construction cohort and do not establish held-out company
generalization. A generalization claim requires a company-disjoint evaluation split fixed before
operator construction.

\subsection{Intervention and Controls}

For an operator $v[p]$, we add a scaled direction to the instruction-span residual states during
prefill,

\begin{equation}
h_{15,p} \leftarrow h_{15,p} + \alpha v[p].
\end{equation}

The transform is disabled during one-token decode steps, so the intervention acts on the prompt
representation rather than repeatedly modifying the generated continuation. The operator study
compares a scalar intervention on one neuron, the one-dimensional DIM direction, and the
four-dimensional cone. The final comparison must use a common dose convention or matched per-token
intervention norm; raw alpha values alone are not comparable across operators.

The controls are matched-norm random directions and identity-stripped prompts. Random controls test
whether a learned direction has a larger effect than an arbitrary residual perturbation.
Identity-stripped prompts test whether the learned direction acts only on a company identifier or
also shifts general investment stance. These controls do not by themselves establish semantic
independence or entity debiasing.

\subsection{Analysis Protocol}

We report margin displacement, monotonicity over the dose sweep, greedy decision flips, JSON parse
rate, and the number of evaluated companies for every arm. Layer localization is selected before the
operator comparison and is evaluated with the same prompt and scoring protocol. Evidence sensitivity
is assessed by paired positive, negative, mixed, and zero-evidence conditions; generated reasons are
qualitative examples rather than a standalone measure of explanation faithfulness.

The primary behavioral claim is based on generated decision flips. Margin movement without a
corresponding greedy decision change is reported as a readout effect and is not counted as a decision
flip.

\section{Layer-Localized Steering Analysis}

We will report a complete span-by-layer residual-patching map for entity, evidence, instruction, and
answer-prefix spans. The steering layer is selected with a rule fixed before the operator comparison,
using fixed-token margin displacement as the primary localization metric and greedy generation as a
secondary check.

The existing Qwen3.5-4B development characterization makes L15 a candidate site: the stance direction
explains $R^2=0.605$ of final margin variation in a 16-company layer scan, with L15 the highest value
among the tested layers. This observation is not a substitute for the complete span-by-layer map,
and the final paper should not claim that the selected layer is the origin of entity-dependent
behavior.

\section{Representation-Steering Operators}

We will compare the single-neuron scalar intervention, DIM, and four-dimensional cone on the same
target set and dose convention. For each arm, we will report the complete dose-response curve,
intervention norm, margin displacement, greedy decision, parse rate, and matched-norm random control.

The current frozen DIM pilot provides a decision-level feasibility result: MO, CNC, and FOXA move
from clean \texttt{sell} outputs toward \texttt{buy}, with observed flips at alpha 4, 5, and 4,
respectively. The same pilot contains one random-direction control and one identity-stripped prompt.
These remain pilot observations until the multi-seed and multi-target confirmation run is stored as
a compact result artifact.

The cone construction artifact is available, but the current persisted cone evaluation is only a
smoke result at alpha 0. Claims about smoother scaling, reduced saturation, or an advantage over a
one-dimensional operator require a new output artifact with dimension ablations and matched-norm
controls. Generated reasons may illustrate different outputs, but they do not identify individual
cone axes with distinct financial concepts.

\section{Decision-Level Steering Boundaries}

This section will organize results around the distinction between fixed-token margin movement and
generated decisions. For every evidence condition, we will report the clean and steered margin, the
generated JSON decision, the flip indicator, and parse status.

The first boundary is readout-to-generation divergence: a positive margin can coexist with a greedy
\texttt{sell}, as in the current DIM pilot for CNC. The second boundary is evidence sensitivity:
strong adverse evidence may delay or prevent a greedy flip even when the intervention moves the
margin in the opposite direction. This second result must be rerun with one frozen renderer because
the current note mixes frozen balanced prompts with custom scenario prompts.

Zero-evidence prompts and identity-stripped prompts can test whether the operator shifts a general
stance under missing company evidence. Macro-context probes are exploratory extensions and should
remain outside the primary claim unless they receive the same prompt, split, and compact-artifact
provenance as the main evaluation.

\section{Discussion}

The central interpretation of this study is that investment decision steering can be analyzed as a
layer-localized representation-control problem. The entity-dependent behavior that motivates the
study does not by itself establish where the behavior originates. Instead, layer-wise localization
identifies a practical intervention site, and operator comparisons test how much control can be
obtained there.

The distinction between fixed-token readout and generated decision is central to the interpretation.
A direction may move the buy/sell margin without changing the greedy JSON output, especially under
strong evidence or formatting constraints. Conversely, a generated flip is a stronger behavioral
outcome than margin movement alone and should be the basis for claims about decision control.

\section{Limitations}

The current study uses one model, a fixed decision template, small construction-cohort pilots, a
limited random-control arm, and greedy decoding. The present operator construction and pilot
evaluation do not establish held-out company generalization. A complete operator comparison still
requires a full-cohort sweep, dimension ablation, matched-norm controls, and cross-model replication.
The work studies interpretability and controllability of an investment decision protocol; it does not
establish entity-selective debiasing, a trading strategy, or financial advice.

\section{Ethical Considerations}

Inference-time steering could change an investment recommendation while leaving a plausible surface
rationale in the generated response. This creates a risk that users may mistake a controlled output
for an independent financial assessment. We therefore treat generated decision flips and rationales
as outputs of an interpretability and controllability experiment, not as evidence of a trading
strategy or financial advice.

\section{Conclusion}

We study where investment decision steering is accessible inside an LLM and how different residual
representation operators control it. The paper combines layer-and-span localization with a comparison
of a single-neuron intervention, a DIM direction, and a multi-dimensional concept cone. By reporting
fixed-token margin movement separately from greedy decision flips, it characterizes both the
availability of an internal steering handle and the limits of translating that handle into generated
behavior.

\appendix

\section{Prompt and Measurement Details}

The appendix will include the canonical prompt format, scenario-template distinction, token-span
definition, margin measurement, greedy parsing rule, and intervention timing.

\section{Reproducibility and Supplementary Materials}

The appendix will describe the anonymized code and data supplement, compact derived artifacts,
operator provenance, and analyses that remain outside the primary paper claims.

\bibliography{custom}

\end{document}
